\section{Discussion}
\label{header:discussion}

\paragraph{}The empirical results of Section~\ref{header:results} substantiate the central hypothesis stated in Section~\ref{sect:meth4}: that a $(t, n)$ threshold-signature protocol grounded in BLS12-381 can deliver the security guarantees of multi-party authorisation without sacrificing the latency profile that users have come to expect from QR-based mobile payments. This section interprets the four streams of evidence --- functional correctness, performance, adversarial resilience, and user acceptance --- and places them in dialogue with the literature reviewed in Section~\ref{header:rrl}. It is organised around five axes of interpretation, each of which corresponds to a thread that recurs across the reviewed gaps in the literature.

\subsection{Cryptographic Correctness vs.\ Operational Security}
\label{sect:disc_crypto}

\paragraph{}A useful way to read the functional and adversarial results in tandem is to separate \textit{cryptographic} guarantees, which are unconditional and follow from the security of BLS and the hardness of the elliptic-curve discrete-log problem, from \textit{operational} guarantees, which depend on how the protocol is wired into the surrounding software and on the trust assumptions placed on the proxy and the persistence layer. The 35-test functional suite (Section~\ref{sect:results_func}) and the offline cryptographic attack group A1--A7 (Section~\ref{sect:results_mitm}) jointly establish the cryptographic guarantees: every legitimate 3-of-5 aggregate verifies, no 2-of-5 partial aggregate verifies, no payload mutation goes undetected, and no off-curve point is admitted by the verifier. These results are deterministic --- they hold for every input the test suite supplies, and would continue to hold under any extension of that suite that respects the protocol's stated message format.

\paragraph{}The online attack group B1--B6 instead exercises operational guarantees. The fact that the verifier rejects a forged \texttt{amount} field (B2, B3) is not a property of BLS as a primitive; it is a property of the patched \texttt{socketHandler.js}, which now reads the authoritative amount from the persisted \texttt{Transaction} document rather than from the client envelope. The original implementation, in which the client-supplied \texttt{amount} was used as the operand of the deduction, would have passed every functional test and every offline cryptographic test, and would still have admitted a non-trivial economic attack against any deployment that lacked transport-layer encryption. This separation, between cryptographic and operational guarantees, is the principal pedagogical takeaway of the present study: a threshold-signature protocol is necessary but not sufficient for an end-to-end secure payment system, because the integrity of the bound transaction \textit{message} depends on the surrounding application logic preserving the field whose value is being authorised. Tests B2 and B3 now provide regression coverage for that property.

\paragraph{}A complementary observation concerns the proxy. By construction, the proxy in this study learns nothing of the partial signatures it relays beyond their existence and timing, since the inter-client traffic is encrypted with AES-256-GCM keyed by an ephemeral Diffie--Hellman exchange (Section~\ref{sect:meth2_b}, item~h). The proxy is therefore consistent with the ``honest-but-curious'' adversary model used in much of the proxy re-encryption literature~\cite{chen2011, ruffing2014}: it can deny service and observe metadata, but cannot read or forge cryptographic content. The audit metadata it persists (session identifiers, signature counts, timestamps) is exactly the residue that the threat model permits.

\subsection{Performance Margins and the 200~ms Acceptance Criterion}
\label{sect:disc_perf}

\paragraph{}The latency benchmark substantially exceeds the $< 200$~ms acceptance criterion. The mean verification latency of $4.83$~ms and the 99th-percentile worst case of $8$~ms leave a margin of more than an order of magnitude against the criterion, even on the localhost build. Two consequences of this margin deserve interpretation.

\paragraph{}First, the protocol leaves substantial headroom for additional cryptographic work without breaching the user-perceptible latency boundary. The verification step is cheap because it is dominated by three near-constant-cost operations: a single \texttt{findOne} lookup on the \texttt{Transaction} collection, a single \texttt{save} that updates the terminal status, and a single broadcast to the room. If, in a future iteration, the verifier were extended to perform additional cryptographic checks at scan time --- for instance, a freshness proof bound to the merchant's session, a verifiable random function-based receipt, or a per-scan zero-knowledge attestation --- the cumulative cost of those additions could grow by tens of milliseconds before the 200~ms criterion would come under threat. The protocol's performance budget is therefore not pinched; it is comfortably absorbing.

\paragraph{}Second, the QR-generation cost (mean $24.73$~ms, p99 $32$~ms) is dominated by the cryptographic primitives themselves rather than by the persistence layer. This is empirically inferred from the relative sizes of the verification and generation latencies: verification involves the same number of MongoDB writes as generation, yet costs roughly five times less, because it does not perform a partial signing, an aggregation, or a self-verification pairing. The implication is that the principal lever for further latency reduction on the generation path is the cryptographic library, not the database. The \texttt{@noble/curves} library used in this study is a pure-JavaScript implementation of BLS12-381; a native (e.g., \texttt{node-gyp}-bound) implementation, or a WebAssembly port, would be expected to reduce QR-generation latency by a meaningful factor, although given the existing margin against the criterion the engineering case for that work is weaker than it would otherwise be.

\paragraph{}A third observation, anticipating Section~\ref{sect:disc_qr}, concerns the additivity of latency. The total end-to-end latency observed in the benchmark (mean $29.57$~ms, p99 $36$~ms) accounts only for the \textit{computational and network} components of the round trip; it does not include the camera autofocus latency that Cycle 1 of the UAT exposed in the Big QR encoding. A 2--3 second autofocus delay, when added to a $\sim 30$~ms backend round trip, dominates the user's experience by two orders of magnitude. The dominant latency in the system, in other words, is not cryptographic --- it is optical.

\subsection{The Security--Scannability Trade-off Surfaced by Iterative UAT}
\label{sect:disc_qr}

\paragraph{}The most operationally consequential finding of the UAT pipeline is the asymmetry in QR recognition latency between the two cycles: three of five Cycle-1 participants (Big QR, full in-band payload) reported a 2--3 second focus delay, while no Cycle-2 participant (Small QR, reference-token payload) did (Section~\ref{sect:results_uat2}). The framing in Sections~\ref{sect:results_uat2} and~\ref{sect:disc_perf} treats this as a security--scannability trade-off, and it is useful here to make the structure of that trade-off explicit, because the protocol design admits both endpoints simultaneously rather than forcing a choice between them.

\paragraph{}The two variants compared across the cycles --- rendered side-by-side at equal physical scale in Figure~\ref{fig:qr-comparison-results} of Section~\ref{sect:results_refine} --- sit at opposite ends of an encoding axis. The Big QR places the entire authorisation envelope --- $\{PK, g, F_{\text{sig}}, M\}$ --- inside the QR matrix itself. This makes the variant \textit{self-contained}: a verifier holding only the QR can, in principle, verify the authorisation without contacting any third party, and there is no separate retrieval channel that an adversary could compromise to substitute a forged payload. The cost is the data density required to encode the payload. BLS12-381 group elements are compact (a $G_1$ point in compressed form is 48 bytes), and the joint public key, generator, and signature together with the transaction message $M$ bring the payload to roughly $440$ bytes once serialised in JSON with field labels; this typically forces the QR to a higher version (denser module grid) than a typical payment QR carrying only a URL or a short token~\cite{worldbank2021ch5, corpuz2025}. The denser the grid, the more precise the camera autofocus must be to discriminate adjacent modules, and the longer the autofocus controller takes to converge --- which is precisely what the 2--3 second focus delay reported by 60\% of Cycle-1 participants reflects.

\paragraph{}The Small QR, by contrast, embeds only a short reference token (in the refined prototype, a UUID and a session identifier together fitting comfortably in a version-2 QR with high error correction). The verifier then retrieves the full authorisation envelope from the proxy at scan time, over the same authenticated channel that the rest of the protocol already trusts. This drops the data density by roughly an order of magnitude, which eliminated the autofocus penalty in Cycle 2 (5/5 participants reported instantaneous recognition). The cost is that the variant is no longer self-contained: a compromised retrieval channel could in principle substitute a forged envelope. In practice, the retrieval channel rides on the same TLS-protected proxy that the rest of the protocol already trusts, so the additional surface is small; nevertheless, it is real, and it shifts a portion of the trust budget from cryptographic content to transport.

\paragraph{}It is also worth flagging the epistemic status of the cross-cycle comparison. Because the Small QR encoding was introduced as a Phase-4 refinement specifically in response to Cycle-1 feedback (Section~\ref{sect:results_refine}), the contrast is not a pre-registered between-subjects A/B in the experimental-design sense; it is an iterative within-pipeline comparison whose treatment assignment was determined by the qualitative signal from the first cohort. The direction of the finding (denser matrix $\rightarrow$ longer autofocus) is robust because it is predicted by the underlying optics of camera autofocus on high-density patterns~\cite{corpuz2025}; the precise proportions, however, should be read as a confirmation of an expected effect rather than as a population estimate.

\paragraph{}The protocol design does not require a global choice between these two endpoints. Because the threshold $t$ is already a per-transaction policy decision (Section~\ref{sect:meth2_b}, item~d), it is straightforward to make the QR variant a similar per-transaction decision keyed by transaction value. This is the recommendation developed in Section~\ref{sect:rec_adaptive}.

\paragraph{}A final remark on Scenario~6 (Expired QR Code) is in order. The observer note --- that the participant ``just need[ed] to open the QR \textit{agad}'' --- conflates two latencies that are otherwise independent: the formal two-minute QR TTL and the 2--3 second autofocus delay. From the participant's perspective, however, the conflation is rational, because the latencies compose: every second spent waiting for the camera to focus is a second consumed from the TTL budget. The autofocus delay therefore functions as a soft tightening of the effective expiry window, even though the protocol's nominal TTL is unchanged. This is one of several places where a usability signal in the UAT carries a load-bearing implication for the security model: a longer autofocus delay makes legitimate users more likely to retry, which in turn places more pressure on the system's replay defences. The replay defences in the current implementation are robust (tests B5 and~B6), but the indirect coupling is worth naming.

\subsection{Perceived Security and the Multi-Party Mental Model}
\label{sect:disc_perceived}

\paragraph{}A consistent finding in the perceived-security literature is that users underweight cryptographic mechanisms that they do not directly observe and overweight procedural mechanisms that they do~\cite{krombholz2014}. The UAT in the present study is consistent with that pattern: every participant rated the multi-party approval workflow as ``significantly safer'' or ``somewhat safer'' than a single-user payment, but no participant referred to BLS, threshold signatures, or pairing-based verification when asked to articulate their reasoning. The mechanism that participants saw --- a pop-up on each approver's phone, a visible quorum, a refusal to proceed without it --- dominated the mechanism they did not see (the cryptographic enforcement that made the visible mechanism load-bearing).

\paragraph{}This has two consequences for the design. First, the visible enforcement of the quorum is a security-critical UI element in its own right; degrading it (for instance, by allowing approvals to be silently auto-confirmed under any circumstance, or by hiding the approver list from the shopper) would be expected to materially reduce perceived security even if the underlying protocol were unchanged. Second, the open-ended feedback's request for stricter credential confirmation at the point of approval (Section~\ref{sect:results_uat_open}) is consistent with this pattern: participants want a visible per-approval check that ratifies the cryptographic step they cannot see. Recommendation~\ref{sect:rec_mfa} is therefore as much a perceived-security recommendation as it is a credential-binding one.

\paragraph{}It is worth noting, however, the limit of this finding. None of the participants in this UAT had cryptographic training; the perceived-security ratings should be read as evidence about lay users' mental models, not about the security properties of the protocol itself. The protocol's security claims rest on the mathematical and adversarial results of Sections~\ref{sect:results_func} and~\ref{sect:results_mitm}, not on user-reported confidence.

\subsection{Limitations and Threats to Validity}
\label{sect:disc_limits}

\paragraph{}Three principal limitations should be flagged for the integrity of the foregoing claims.

\paragraph{a.) Sample size, cohort independence, and instrument adaptation.}
The UAT was conducted across two sequential cycles of $n = 5$ participants each, for a total of ten distinct participants. This is sufficient for a formative usability evaluation in the tradition of small-sample SUS studies~\cite{brooke1996sus, bangor2008sus}, but it is not sufficient to support inferential claims about population-level usability or perceived security. The cross-cycle asymmetry observed in the QR-recognition item (Section~\ref{sect:results_uat2}) is large enough that the qualitative direction of the finding is unlikely to reverse under a larger sample, but the precise proportions should not be taken as point estimates of any underlying rate. The cross-cycle contrast is also not a pre-registered between-subjects A/B: the encoding change was introduced as a Phase-4 refinement in response to Cycle-1 feedback (Section~\ref{sect:disc_qr}), so the comparison should be read as a confirmatory check that an expected effect was observed, rather than as an independent test of two pre-existing conditions. The six-item SUS adaptation, in particular, departs from the canonical 10-item instrument, and the resulting score is therefore not strictly interchangeable with normative SUS benchmarks; it is reported in Section~\ref{sect:results_uat_sus} with that caveat.

\paragraph{b.) Latency benchmark scope.}
The latency benchmark was conducted on localhost, eliminating network round-trip time as a variable. The production deployment introduces an additional 50--150~ms of geographic routing overhead, which the discussion in Section~\ref{sect:disc_perf} addresses qualitatively. A controlled benchmark on the production environment, or on a representative cloud-to-mobile path, would be a useful confirmation. In addition, the benchmark exercises the low-value (threshold $t = 1$) path; the high-value path ($t \geq 2$) introduces additional inter-client coordination, whose latency profile is bounded above by the slowest approver's response time and is therefore dominated by human factors rather than by cryptography.

\paragraph{c.) Threat model boundary.}
The MITM evaluation models a network-level adversary who has either intercepted plaintext (in non-TLS deployments) or broken TLS (in production deployments), and a cryptographic adversary who can submit arbitrary candidate signatures and message digests to the verifier. It does not model an endpoint-compromise adversary who has obtained one of the five private key shares from a participant device; in that scenario, the threshold property continues to provide protection (a single share is insufficient to forge), but a coalition of $t$ compromised devices would, by construction, be able to authorise transactions. Defences against the endpoint-compromise scenario --- proactive secret sharing, hardware-backed key storage, behavioural anomaly detection on signing patterns --- are out of scope for the present study and are noted in Section~\ref{sect:rec_proactive} and Section~\ref{sect:rec_hardware} as directions for follow-up work.
